\chapter{Process Flow: Surprise Gating and Ranking (``Forgetting'')}

\section{The Attention and Surprise Gate}
\textbf{Trigger:} Any inbound turn, before it reaches the Observer.
\textbf{Action / UI State:} The Panel does \emph{not} render routine turns; a suppressed turn shows a
dimmed ``surprise 0.12 $\rightarrow$ skipped'' chip, while a surprising one lights up.
\textbf{System Logic:} The \textbf{Attention Gate} discards noise (tool heartbeats, acknowledgements),
and the \textbf{Surprise Gate} computes a prediction error --- the cosine distance between the turn's
embedding and the observer's expectation. Turns below \texttt{SURPRISE\_FLOOR} are not encoded.
\textbf{Data Flow \& Mutators:} A suppressed turn produces no \texttt{memories} row; a
\texttt{memory\_event} records the decision so the panel and the eval ledger can show \emph{why}
nothing was stored.

\section{A Surprising Turn Is Encoded}
\textbf{Trigger:} A turn whose surprise clears the floor.
\textbf{Action / UI State:} The Panel shows the surprise score alongside the Observer verdict.
\textbf{System Logic:} Surprise enters the salience equation as a multiplier, so an unexpected fact
outranks a routine one at equal persistence and consequence. A user \emph{correction} triggers the
reward-prediction-error path and is consolidated immediately.
\textbf{Data Flow \& Mutators:} \texttt{memories.surprise} is set; the row enters the write path with
elevated salience.

\section{Ranking, Not Deleting}
\textbf{Trigger:} The Ranker runs (on consolidation, or periodically).
\textbf{Action / UI State:} The Panel reports the ranking summary (``4 memories dropped out of the
active set''), and each can be clicked to prove it is \emph{still there}.
\textbf{System Logic:} Each memory's \texttt{utility} is recomputed:
\begin{Verbatim}[breaklines=true,fontsize=\small]
utility = exp(-ln2 * age / half_life) * (1 + log(1 + access_count))
# low utility -> lower retrieval rank + leave the attention window (STILL AVAILABLE)
\end{Verbatim}
\emph{Crucially, nothing is destroyed.} A deprioritised memory is available but not accessible; the
moment a matching cue arrives, the Activator revives it and it surfaces again. Hard deletion happens
only on an explicit right-to-be-forgotten request.
\textbf{Data Flow \& Mutators:} \texttt{memories.utility} updated; the memory's \emph{rank} changes,
not its existence.

\section{The Cue Always Works}
\textbf{Trigger:} A later turn whose cue overlaps a deprioritised memory.
\textbf{Action / UI State:} The Panel shows the memory returning with a ``revived (utility 0.08
$\rightarrow$ served)'' note.
\textbf{System Logic:} This is the deliberate inversion of the human tip-of-the-tongue failure. Human
memory fails because a weak cue cannot cross the access threshold; Tottie's search is exhaustive, so
the right cue always retrieves. The memory was never lost --- it was waiting.
\textbf{Data Flow \& Mutators:} A read; \texttt{access\_count} increments, which raises utility and
promotes the memory back into the active set (the testing effect).

\section{Interference Control}
\textbf{Trigger:} Two memories with overlapping embeddings.
\textbf{System Logic:} On write, \textbf{sparse pattern separation} decorrelates near-identical
representations; on read, category/namespace filters keep unrelated domains from polluting each other.
This is the artificial analogue of proactive and retroactive interference control --- and unlike the
brain, it is a deliberate, tunable design choice.
\textbf{Data Flow \& Mutators:} A separation key or namespace is recorded on the row; conflicting
rows are merged or superseded rather than stacked.
